\subsection{Stenotic Hemodynamics and Relevant Observables}
\label{subsec:stenotic-hemodynamics-observables}

Stenosis is a geometric and hemodynamic problem: a local lumen reduction can
alter axial velocity, pressure drop, wall loading, recirculation, and the
relationship between flow and downstream perfusion. Clinical and computational
literatures therefore use different observables depending on the question.
Pressure-ratio work emphasizes pressure drops and fractional-flow style
indices
\cite{PijlsEtAl1996FFRFunctionalSeverity,ToninoEtAl2010FAMEAngioFunctional,TebaldiEtAl2015FFROverview},
whereas computational fluid-dynamics studies can expose local velocity,
traction, and wall-shear structures that are not section-averaged quantities
\cite{JohnEtAl2017WallShearStressVectorForm,GaldiEtAl2008HemodynamicalFlows}.
The present report uses velocity observables only. It does not report pressure
drop, FFR, CT-FFR, wall shear stress, or clinical validation.

Reduced models are attractive because stenosis studies often require repeated
parameter sweeps, reproducible solver comparisons, or network-level coupling.
The cost is that the relevant output must be compatible with the reduced state.
For a one-dimensional area--flow model, section mean velocity and physical flow
are natural; resolved radial profile, secondary flow, skewness, and wall
traction are not native outputs unless an additional reconstruction or
observation map is declared.

\subsection{State-of-the-Art Model Families}
\label{subsec:model-hierarchy-literature}

The current modeling landscape is best read as a hierarchy of model roles, not
as a single ladder of accuracy. Three-dimensional CFD and FSI models retain
geometry-resolved velocity and stress fields and can represent local stenotic
acceleration, separation, and wall interaction more directly than reduced
models
\cite{CourchaineRugonyi2018CardiacDevelopmentCFD,DaveyEtAl2024CardiacFluidDynamicsHeart}.
They also require detailed geometry, material assumptions, boundary histories,
mesh generation, and substantially higher computational cost. Coupled 3D--1D
or geometrical multiscale formulations address domain-size and boundary
limitations by joining local resolved regions with reduced vascular networks
\cite{FormaggiaEtAl2007FSICoupling,QuarteroniVeneziani2003GeometricalMultiscale,KopplHelmig2023DimensionReducedModeling}.

One-dimensional distributed models retain axial propagation and
pressure--area--flow dynamics while replacing the transverse velocity field by
closure terms
\cite{FormaggiaEtAl2003OneDimensionalBloodFlow,QuarteroniFormaggia2004CardiovascularModeling,ShiEtAl2011BloodFlowModelReview}.
Zero-dimensional networks and terminal models are still cheaper and are useful
for lumped impedance and perfusion effects, but they do not resolve axial
stenosis localization. Intermediate 2D or axisymmetric models can retain more
local structure than 1D systems with less cost than full 3D, but they still
require their own symmetry and boundary assumptions.

Recent surrogate, PINN, and operator-learning methods add another axis:
trained maps can reduce evaluation cost or assimilate data, but their evidence
depends on training data, uncertainty treatment, boundary-condition coverage,
and comparison with mechanistic solvers or measurements
\cite{RaissiEtAl2019PhysicsInformedNeuralNetworks,LuEtAl2021DeepONet,YuEtAl2026PINNHemodynamicsReview,VelikorodnyEtAl2025DeepOperatorStenosed}.
They are therefore useful state-of-the-art context and future-work vocabulary
for stenosed-vessel modeling, but they are not used as forward generators in
the present report.

\begin{table}[!htb]
    \centering
    \scriptsize
    \caption{State-of-the-art model and method families relevant to this
    report. The table maps field context to the evidence standard needed before
    a result can support a claim in this thesis.}
    \label{tab:sota-model-method-map}
    \begin{tabularx}{\textwidth}{@{}
        >{\raggedright\arraybackslash}p{0.21\textwidth}
        >{\raggedright\arraybackslash}p{0.27\textwidth}
        >{\raggedright\arraybackslash}p{0.25\textwidth}
        >{\raggedright\arraybackslash}X@{}}
        \toprule
        \textbf{Family}
            & \textbf{Representative role}
            & \textbf{Evidence requirement}
            & \textbf{Reading in this report} \\
        \midrule
        3D CFD / FSI and multiscale 3D--1D--0D
            & Resolved local fields, wall interaction, and coupled circulation
            \cite{FormaggiaEtAl2007FSICoupling,DaveyEtAl2024CardiacFluidDynamicsHeart}.
            & Matched geometry, wall and material data, boundary histories,
            mesh/time studies, and output operators.
            & Reference vocabulary and comparison target, not the selected
            repeated forward generator. \\
        Distributed 1D stenosis-aware models
            & Low-cost area--flow dynamics with closure-indexed wall, profile,
            friction, and stenosis terms
            \cite{FormaggiaEtAl2003OneDimensionalBloodFlow,CanicEtAl2024Extended1DStenosis}.
            & Explicit closure contract, boundary rule, coordinate map, and
            verification of the computed realization.
            & Selected model class; the implementation is an
            $R_{\max}$-normalized Canic-derived variant. \\
        Open and reproducible 1D solver workflows
            & Published solver surfaces, benchmark cases, and reusable
            finite-volume infrastructure
            \cite{BoileauEtAl2015Benchmark,BenemeritoEtAl2024OpenBF}.
            & Runnable commands, persisted descriptors, benchmark outputs,
            hashes, and documented optional inputs.
            & Workflow comparator and reproducibility standard, not an
            independent-solver parity claim. \\
        Well-balanced and high-order 1D numerical schemes
            & Equilibrium-preserving finite-volume or high-order discretizations
            for elastic-vessel networks and source terms
            \cite{LuccaEtAl2025FFRPrediction,PimentelGarciaEtAl2025HighOrderWellBalanced}.
            & Rest-state/equilibrium preservation, manufactured solutions,
            refinement, and boundary-regime diagnostics.
            & Context for why the present rest-state drift is a material
            numerical limitation. \\
        Surrogates, PINNs, and operator learning
            & Fast learned maps or physics-constrained reconstructions for
            hemodynamic quantities
            \cite{RaissiEtAl2019PhysicsInformedNeuralNetworks,YuEtAl2026PINNHemodynamicsReview,VelikorodnyEtAl2025DeepOperatorStenosed}.
            & Training-data provenance, uncertainty, out-of-distribution tests,
            and comparison with accepted reference outputs.
            & Bounded outlook only; no learned model is part of the completed
            numerical study. \\
        \bottomrule
    \end{tabularx}
\end{table}

\begin{table}[!htb]
    \centering
    \scriptsize
    \caption{Model-tier roles relevant to the present report. The listed
    observables are admissible outputs only after the corresponding model and
    observation map have been declared.}
    \label{tab:model-tier-review}
    \begin{tabularx}{\textwidth}{@{}
        >{\raggedright\arraybackslash}p{0.19\textwidth}
        >{\raggedright\arraybackslash}p{0.22\textwidth}
        >{\raggedright\arraybackslash}p{0.22\textwidth}
        >{\raggedright\arraybackslash}X@{}}
        \toprule
        \textbf{Tier} & \textbf{Retained physics} & \textbf{Data and cost} & \textbf{Typical admissible outputs} \\
        \midrule
        3D CFD / FSI
            & Resolved velocity, local gradients, geometry, and possibly wall motion.
            & Detailed mesh, material and boundary histories; high cost.
            & Local velocity, pressure, traction, wall shear stress, section averages. \\
        Axisymmetric / 2D
            & More local structure than 1D under symmetry or planar assumptions.
            & Intermediate geometry and mesh burden.
            & Axial/radial fields consistent with the symmetry assumption. \\
        1D distributed
            & Axial wave propagation, area--flow dynamics, wall and friction closures.
            & Vessel centerline, radius, wall law, terminal and inlet data; low cost.
            & Area, flow, section mean velocity, reduced pressure variables. \\
        0D network
            & Lumped resistance, compliance, inertance, and terminal coupling.
            & Lowest spatial data burden; low cost.
            & Compartment pressures and flows, not local stenotic fields. \\
        Geometrical multiscale
            & Local resolved region coupled to reduced upstream/downstream networks.
            & Requires interface conditions and compatible observables.
            & Mixed local and network outputs with interface consistency checks. \\
        \bottomrule
    \end{tabularx}
\end{table}

\subsection{One-Dimensional Models for Nonuniform and Stenotic Vessels}
\label{subsec:literature-1d-stenotic-vessels}

The defining step in a 1D vessel model is cross-sectional averaging. The state
typically stores a section area and volumetric flow, while the discarded
transverse structure re-enters through closure choices: momentum-flux factor,
velocity profile, wall law, friction, rheology, and source terms for nonuniform
reference geometry
\cite{FormaggiaEtAl2003OneDimensionalBloodFlow,ShiEtAl2011BloodFlowModelReview}.
These choices are not interchangeable. A profile factor changes the convective
flux; a wall law changes the pressure--area relation and wave speed; a
friction or rheology law changes dissipative closure; and a stenosis-loss or
variable-radius term changes how geometry enters the reduced momentum balance.

Classical stenosis literature also shows that reduced pressure-loss formulas
and unsteady stenosis corrections can be useful but model a different closure
object than a distributed 1D finite-volume system
\cite{YoungTsai1973SteadyStenosis,YoungTsai1973UnsteadyStenosis,LyrasLee2021PressureDropROM}.
Recent stenosis-focused reduced models, including the Canic et al. extended 1D
formulation, add variable-radius contributions intended to better represent
stenotic geometry in a distributed balance
\cite{CanicEtAl2024Extended1DStenosis}. Coronary and arterial-tree models make
the same point at network scale: useful predictions depend on the declared
wall, terminal, and closure assumptions
\cite{DuanmuEtAl2019CoronaryArterialTree,Wang2014BloodFlowNetworks}.

\begin{table}[!htb]
    \centering
    \scriptsize
    \caption{Closure dimensions that should be separated when comparing
    stenosis-aware 1D models.}
    \label{tab:one-dimensional-closure-review}
    \begin{tabularx}{\textwidth}{@{}
        >{\raggedright\arraybackslash}p{0.20\textwidth}
        >{\raggedright\arraybackslash}p{0.30\textwidth}
        >{\raggedright\arraybackslash}X@{}}
        \toprule
        \textbf{Closure dimension} & \textbf{Common choices} & \textbf{Effect on evidence} \\
        \midrule
        State variables
            & Area--flow, radius--flow, pressure--flow, or network compartment variables.
            & Determines which outputs are native and which require reconstruction. \\
        Wall model
            & Fixed geometry, algebraic pressure--area relation, FSI coupling, terminal compliance.
            & Controls wave speed and pressure interpretation. \\
        Profile / momentum flux
            & Flat, parabolic, power-law, empirical correction factors.
            & Controls convective flux and any reconstructed radial profile. \\
        Friction / rheology
            & Newtonian, Carreau, Carreau--Yasuda, Casson, power-law, tube-flow laws.
            & Controls dissipation and shear-rate proxy interpretation
            \cite{PriesEtAl1992BloodViscosityTubeFlow,KumarEtAl2019CarreauYasudaHemodynamics,LiuEtAl2021ICASRheology}. \\
        Stenosis treatment
            & Reference-radius variation, empirical loss terms, variable-radius distributed terms.
            & Determines whether geometry appears only in area or also as a momentum source. \\
        Numerical method
            & Finite volume, discontinuous Galerkin, WENO, Lax--Friedrichs-type,
            well-balanced, and method-specific boundaries.
            & Determines conservation, equilibrium preservation, stability evidence,
            and boundary evidence needs
            \cite{LeVeque2002FiniteVolumeMethodsForHyperbolicProblems,BeckersKolbe2025LaxFriedrichsHemodynamics,LuccaEtAl2025FFRPrediction}. \\
        Validation target
            & Benchmarks, manufactured solutions, 3D fields, pressure ratios, clinical measurements.
            & Determines whether the result is verification, cross-model comparison, or validation. \\
        \bottomrule
    \end{tabularx}
\end{table}
